
      \documentclass[fleqn]{article}      
      \usepackage{amsmath}
      \usepackage{fontspec}
      \usepackage{kotex} % 한국어 지원  

      \begin{document}      
      쉬운 문제: \\\\  
1. 행렬의 정의는 무엇인가요? \\\\
2. 2x2 행렬의 형태는 어떻게 되나요? \\\\
3. 행렬의 덧셈이 가능한 조건은 무엇인가요? \\\\
4. 행렬의 곱셈에서 교환 법칙이 성립하는가요? \\\\
5. 주대각선의 원소들이 모두 0이 아닌 행렬을 무엇이라고 하나요? \\\\

중간 문제: \\\\  
1. 행렬 A와 B가 같다는 것은 어떤 조건을 만족해야 하나요? \\\\
2. 3x3 행렬의 행렬식이 0이면 어떤 의미인가요? \\\\
3. 행렬 A와 B의 곱 AB를 계산하기 위해 필요한 조건은 무엇인가요? \\\\
4. 전치행렬의 정의는 무엇인가요? \\\\
5. 주어진 행렬의 역행렬을 구하는 방법은 무엇인가요? \\\\

어려운 문제: \\\\  
(어려운 문제는 없습니다) \\\\

객관식 문제: \\\\  
쉬움: \\\\  
1. 행렬의 곱셈에서 결과행렬의 크기는 어떻게 결정되나요? \\\\
   a) 첫 번째 행렬의 행 수와 두 번째 행렬의 열 수 \\\\
   b) 두 행렬의 행 수 \\\\
   c) 두 행렬의 열 수 \\\\
  
2. 다음 중 단위 행렬의 정의는 무엇인가요? \\\\
   a) 모든 원소가 0인 행렬 \\\\
   b) 주대각선의 원소가 1이고 나머지는 0인 행렬 \\\\
   c) 모든 원소가 1인 행렬 \\\\
  
3. 2x2 행렬 A가 다음과 같을 때: A = \(\begin{bmatrix} 1 & 2 \\ 3 & 4 \end{bmatrix}\) 행렬 A의 행렬식은 얼마인가요? \\\\
   a) -2 \\\\
   b) 10 \\\\
   c) 5 \\\\
  
중간: \\\\  
1. 두 행렬 A와 B의 곱이 정의되기 위한 조건은 무엇인가요? \\\\
   a) A의 열 수가 B의 행 수와 같아야 한다. \\\\
   b) A와 B의 크기가 같아야 한다. \\\\
   c) A의 행 수가 B의 열 수와 같아야 한다. \\\\
  
2. 다음 중 전치행렬의 성질이 아닌 것은 무엇인가요? \\\\
   a) (A^T)^T = A \\\\
   b) (AB)^T = B^T A^T \\\\
   c) (A + B)^T = A^T + B \\\\
  
3. 3x3 행렬의 대각합은 무엇인가요? \\\\
   a) 주대각선의 원소들의 합 \\\\
   b) 모든 원소의 합 \\\\
   c) 주대각선의 원소의 곱 \\\\
  
4. 다음 중 역행렬의 정의는 무엇인가요? \\\\
   a) A*A^{-1} = I \\\\
   b) A^{-1}*A = 0 \\\\
   c) A*A^{-1} = 0 \\\\
  
5. 행렬 A의 고유값을 구하기 위한 방법은 무엇인가요? \\\\
   a) A - λI = 0의 해를 구한다. \\\\
   b) 행렬의 전치행렬을 구한다. \\\\
   c) 행렬식을 계산한다. \\\\

      
      \end{document} 
  